\chapter{INTRODUCTION}

The axioms of Geometry were formerly regarded as laws of thought
which an intelligent mind could neither deny nor investigate. Not
only were the axioms to which we have been accustomed found to agree
with our experience, but it was believed that we could not reason on
the supposition that any of them are not true, it has been shown,
however, that it is possible to take a set of axioms, wholly or in
part contradicting those of Euclid, and build up a Geometry as
consistent as his.

We shall give the two most important Non-Euclidean
Geometries.\footnote{See Historical Note, p.~\pageref{histnote}.} In
these the axioms and definitions are taken as in Euclid, with the
exception of those relating to parallel lines. Omitting the axiom on
parallels,\footnote{See p.~\pageref{3hypos}.} we are led to three
hypotheses; one of these establishes the Geometry of Euclid, while
each of the other two gives us a series of propositions both
interesting and useful. Indeed, as long as we can examine but a
limited portion of the universe, it is not possible to prove that
the system of Euclid is true, rather than one of the two
Non-Euclidean Geometries which we are about to describe.

We shall adopt an arrangement which enables us to prove first the
propositions common to the three Geometries, then to produce a
series of propositions and the trigonometrical formulæ for each of
the two Geometries which differ from that of Euclid, and by
analytical methods to derive some of their most striking properties.

We do not propose to investigate directly the foundations of
Geometry, nor even to point out all of the assumptions which have
been made, consciously or unconsciously, in this study. Leaving
undisturbed that which these Geometries have in common, we are free
to fix our attention upon their differences. By a concrete
exposition it may be possible to learn more of the nature of
Geometry than from abstract theory alone.

\newpage
Thus we shall employ most of the terms of Geometry without
repeating the definitions given in our text-books, and assume that
the figures defined by these terms exist. In particular we assume:
\smallskip

\begin{enumerate}
\item[I.] \emph{The existence of straight lines determined by any two
points, and that the shortest path between two points is a straight
line.}

\item[II.] \emph{The existence of planes determined by any three points
not in a straight line, and that a straight line joining any two
points of a plane lies wholly in the plane.}

\item[III.] \emph{That geometrical figures can be moved about without
changing their shape or size.}

\item[IV.] \emph{That a point moving along a line from one position
to another passes through every point of the line between, and that
a geometrical magnitude, for example, an angle, or the length of a
portion of a line, varying from one value to another, passes through
all intermediate values.}
\end{enumerate}

In some of the propositions the proof will be omitted or only the
method of proof suggested, where the details can be supplied from
our common text-books.

